\documentclass[10pt,a4paper]{amsart}
\usepackage[margin=19mm]{geometry}
\usepackage{amsmath,amssymb,mathtools}
\usepackage[scr=rsfs,cal=euler]{mathalfa}
\usepackage{xfrac}
\usepackage[unicode,hidelinks,bookmarks=false]{hyperref}
\newcommand{\ie}{i.e.\ }
\newcommand{\cf}{cf.\ }
\newcommand{\resp}{respectively\ }
\newcommand{\Subsection}{\subsection}
\providecommand{\nobreakdash}{\nobreak\hbox{-}}
\setlength{\emergencystretch}{2em}
\multlinegap0pt
\usepackage{amssymb}
\usepackage{enumerate}
\usepackage{mathtools}
\usepackage[shortlabels]{enumitem}
\usepackage{xcolor}
\newtheorem{theorem}{Theorem}
\newcommand{\Emptyset}{\text{\o}}
\newcommand{\Si}{\Sigma}
\newcommand{\Z}{\mathbb{Z}}
\newcommand{\eqdef }{\overset{\mbox{\tiny{def}}}{=}}
\newcommand{\highG}{\Theta}
\newcommand{\irnp}[2]{\left\langle {#1},{#2} \right\rangle}
\newcommand{\linL}{\mathcal{L}}
\newcommand{\micro}{{\bf \{I- P\}}}
\def\norm#1{\big\Vert#1\big\Vert}
\newcommand{\pZ}{|\pv|}
\newcommand{\pv}{p}
\newcommand{\qq}{\mathfrak{q}}
\newcommand{\qv}{q}
\newcommand{\rth}{{\mathbb{R}^3}}
\newcommand{\tTq}{t^{3\qq}}
\newcommand{\tmTq}{t^{-3\qq}}
\newcommand{\tmq}{t^{-\qq}}
\newcommand{\tq}{t^{\qq}}
\newcommand{\tqq}{t^{2\qq}}
\title{Future global stability of Maxwell--J{\"u}ttner equilibria and vacuum for the massless Boltzmann equation on FLRW spacetimes}
\author{Robert M. Strain$^{\ddagger}$, Martin Taylor$^{\dagger}$, Renato Velozo Ruiz$^{\dagger}$}
\date{Source: arXiv:2606.00175v1 (CC BY 4.0)}
\begin{document}
\maketitle

\noindent\emph{Source and licence.} Future global stability of Maxwell--J{\"u}ttner equilibria and vacuum for the massless Boltzmann equation on FLRW spacetimes. Version arXiv:2606.00175v1 (CC BY 4.0), distributed by arXiv under the Creative Commons Attribution 4.0 International licence, http://creativecommons.org/licenses/by/4.0/, as recorded in the arXiv licence metadata for this exact version and shown on the abstract page https://arxiv.org/abs/2606.00175v1. Full licence notice: \texttt{licenses/arxiv-2606.00175-CC-BY-4.0.txt}.\par\smallskip

\noindent\textit{Editorial scope.} Counted unit: the article's theorem thmmacroscop (source0.tex lines 3048--3057). The article's complete original proof of it is delivered with it (source0.tex lines 3060--3647, one \texttt{proof} environment of 588 lines). No mathematics is written, repaired or re-derived: the statement and the proof are the article's own lines, in the article's own notation, and no retained line is substituted or re-flowed.  Retained uncounted as the source-local context the proof needs: lines 599--607; lines 620--625; lines 1716--1728; lines 1741--1756; lines 1908--1923; lines 2155--2180; lines 2182--2188; lines 2878--2881; lines 2883--2886.  Nothing was omitted from any retained span, so the package carries no omission record.  No retained line contains a citation, so the wrapper carries no bibliography and no bibliography entry.  No retained line contains a figure environment, an image-inclusion command or drawing code, so no image is delivered and the package declares no asset dependency.  Editorial formatting only, in addition: an AMS \texttt{amsart} preamble that restates the article's own declarations and macros used by the retained lines, namely the environments \texttt{theorem} on separate counters, together with the standard packages those lines need.  The retained environments print in this document as Theorem 1 in order of appearance, whereas the article prints the same items with the numbering of its own full document.  The complete native source of the article as distributed in the arXiv e-print for this exact version is bundled unchanged as \texttt{sources/201738/source0.tex}.
\begin{align}
        \partial_t f+ \frac{\pv^i}{|\tq \pv|}\partial_{x^i} f
    -\frac{2\qq }{t} p^i \partial_{p^i} f
    +
    \linL f
=
\Gamma (f,f),
\label{rBoltz00}
\end{align}

\begin{align}
\Gamma (h_1,h_2)
&\eqdef
 J^{-1/2}Q(\sqrt{J} h_1,\sqrt{J} h_2).
\label{gamma0}
\end{align}

\begin{equation}\label{defmacrocoef}
\begin{split}
    \mathcal{A}^f(t,x) & = \tTq\int_\rth f(t,x,\pv)\sqrt{J(\tqq p)}dp,
    \\
    \mathcal{B}_i^f (t,x)
    & =
    \frac{\tTq}{2}\int_\rth f(t,x,\pv) \tqq p^i\sqrt{J(\tqq p)}dp, 
    \\
    \mathcal{C}^f(t,x)
    &=
    \frac{\tTq}{\sqrt{3}}\int_{\rth}f(t,x,\pv)( |\tqq \pv |-3)\sqrt{J(\tqq p)} dp.
\end{split}
\end{equation}

\begin{equation}\label{orthogonal.basis.use}
    \begin{split}
         \tTq \mathfrak{e}_{j+5}(\tqq \pv) \eqdef &~ 
   \sqrt{\frac{3}{7}}
\left(\frac{p^j}{{\pZ}} -\tqq \pv^j\right)\tTq\sqrt{J(\tqq p)},\qquad  \text{for }\quad 1\leq j \leq 3,
      \\ 
   \tTq \mathfrak{e}_{9}(\tqq \pv) \eqdef &~ 
{t^{2\qq}} \frac{\sqrt{5}}{2}\left(\frac{(p^1)^2 -(p^3)^2}{{|p|}} \right)t^{3\qq}\sqrt{J(t^{2\qq} p)},
      \\ 
   \tTq \mathfrak{e}_{10}(\tqq \pv) \eqdef &~ 
{t^{2\qq}} \frac{\sqrt{5}}{2} \left(\frac{(p^1)^2+(p^3)^2 -2(p^2)^2}{{|p|}} \right)t^{3\qq}\sqrt{J(t^{2\qq} p)},
      \\ 
   \tTq \mathfrak{e}_{i+j+8}(\tqq \pv) \eqdef &~ 
\frac{\sqrt{5}}{2}{\tqq}\frac{ \pv^l \pv^{j}}{|p|}\tTq\sqrt{J(\tqq p)},\qquad  \text{for }\quad  1\leq l<j \leq 3.
    \end{split}
\end{equation}

\begin{align} \notag
\sqrt{\frac{7}{3}}\tmq\left(\partial_j\mathcal{A}
-
\sqrt{3}\partial_j\mathcal{C} \right)
&=-\tmTq  \partial_t\left( \tTq \mathfrak{m}_{j+5} \right)+\mathfrak{u}_{j+5}, 
\\  \label{macroscopic}
\tmq\frac{1}{2\sqrt{5}} \left( \partial_1 \mathcal{B}_1 - \partial_3 \mathcal{B}_3\right)
&=-\tmTq  \partial_t\left( \tTq \mathfrak{m}_{9} \right)+\mathfrak{u}_{9},
 \\ \notag
\tmq\frac{1}{6\sqrt{5}}\left( \nabla_x\cdot \mathcal{B}-3\partial_2 \mathcal{B}_2\right)
&=-\tmTq  \partial_t\left( \tTq \mathfrak{m}_{10} \right)+\mathfrak{u}_{10},
 \\ \notag
\frac{1}{\sqrt{5}}  \tmq\left(\partial_j \mathcal{B}_{l}+\partial_l \mathcal{B}_{j} \right)
&=-\tmTq  \partial_t\left( \tTq \mathfrak{m}_{l+j+9} \right)+\mathfrak{u}_{l+j+9},\quad \text{for}
\quad (l\neq j),
\end{align}

\begin{equation}\label{local.cons.alt}
\begin{split}
 \partial_t  \mathcal{A} 
     +
\frac{1}{2}  \tmq \nabla_x \cdot \mathcal{B}
+
\frac{3\qq }{t} {\mathcal{A}}
&=-\tmq \nabla_x \cdot \highG_{1},
\\ 
\partial_t \mathcal{B}_{j} 
 +
2\tmq\partial_j \left( \mathcal{A}
-\frac{2}{\sqrt{3}}
\mathcal{C}  \right)
+
\frac{3\qq }{t} {\mathcal{B}_{j}}
&=-\tmq \nabla_x \cdot \highG_{j+1}, \quad \text{for}\quad 1 \leq j \leq 3,
\\ 
\partial_t  \mathcal{C} 
  +
\frac{1}{2\sqrt{3}}  \tmq \nabla_x \cdot \mathcal{B} 
+
\frac{3\qq }{t} {\mathcal{C}}
&=-\tmq \nabla_x \cdot \highG_{5}.
\end{split}
\end{equation}

\begin{align}\label{macroscopic.B}
 \tmTq\partial_t \left(\tTq \mathcal{B}_{j} \right) 
 +
\tmq\frac{2}{\sqrt{3}}\partial_j \mathcal{C}  
&=
-\tmTq  \partial_t\left( \tTq \mathfrak{m}_{j+13} \right)+\mathfrak{u}_{j+13}\quad  (1 \leq j \leq 3),
\end{align}

\begin{align}\label{dispersion.term.basis.est}
   \left| \left\langle \micro  f, \tTq \mathfrak{e}_{\ell}(\tqq \pv)\right\rangle \right|
        \lesssim &~   \norm{J^{\frac{1}{4}} \micro   f}_{L^2_{\pv}},
\end{align}

    \begin{align}\label{lin.basis.est}
   \left| \left\langle \linL( \micro  f), \tTq \mathfrak{e}_{\ell}(\tqq \pv)\right\rangle \right|
        \lesssim &~ \tmTq \norm{J^{\frac{1}{8}} \micro  f}_{L^2_{\pv}}.
\end{align}

\begin{theorem}[Estimates for macroscopic quantities]\label{thmmacroscop}
Consider \(T>1\), and a suitably regular solution \(f\) of \eqref{rBoltz00} on \([1,T]\times \mathbb{T}^3_x\times \mathbb{R}^3_p\). Then the macroscopic quatities defined by \eqref{defmacrocoef} satisfy the following uniform estimate
\begin{align}\label{abcest.torus1}
\norm{t^{3\qq/2} [\mathcal{A},\mathcal{B},\mathcal{C}]}_{L^1_k L^2_T} &\lesssim
\Vert t^{3\qq/2}\micro f \Vert_{L^1_k L^2_T L^2_{\pv}} +\Vert \tTq f \Vert_{L^1_k L^\infty_T L^2_{\pv}} + \Vert f_1 \Vert_{L^1_k L^2_{\pv}} \\ \notag
&\qquad+ \int_{\mathbb{Z}^3_k} \left( \int^{T}_{1} t^{5\qq}\left| \widehat{\mathfrak{h}}(t,k)\right|^2
 dt \right)^{\frac12} d\Sigma (k),
\end{align}
where \(\norm{ [\mathcal{A},\mathcal{B},\mathcal{C}]}\eqdef\norm{ \mathcal{A}}+\norm{ \mathcal{B}}+\norm{ \mathcal{C}}\), and the definition of $ \widehat{\mathfrak{h}}(t,k)$ is provided below in \eqref{rhs.def.terms.FT}, \eqref{def.GaF}  and \eqref{index.sum.notation} with \eqref{orthogonal.basis.use}. 
\end{theorem}

\begin{proof}
First, we take the Fourier transform in \eqref{macroscopic} and \eqref{macroscopic.B} to obtain 
\begin{align} \label{macroscopic.FT.B}
 \tmTq\partial_t \left(\tTq \widehat{\mathcal{B}}_j \right) 
 +
\frac{2}{\sqrt{3}}\tmq ik_j  \widehat{\mathcal{C}} 
&=
-\tmTq  \partial_t\left( \tTq \widehat{\mathfrak{m}}_{j+13} \right)+\hat{\mathfrak{u}}_{j+13} , \quad (1 \leq j \leq 3),
\\ 
\label{macroscopic.FT.AC}
\tmq\sqrt{\frac{7}{3}}ik_j \left(\widehat{\mathcal{A}}
-
\sqrt{3}\widehat{\mathcal{C}} \right)
&=-\tmTq  \partial_t\left( \tTq \widehat{\mathfrak{m}}_{j+5} \right)+\hat{\mathfrak{u}}_{j+5},
\\ \label{macroscopic.FT.Bthree}
\tmq\frac{i}{2\sqrt{5}} \left( k_1 \widehat{\mathcal{B}}_1 - k_3 \widehat{\mathcal{B}}_3\right)
&=-\tmTq  \partial_t\left( \tTq \widehat{\mathfrak{m}}_{9} \right)+\widehat{\mathfrak{u}}_{9},
 \\ \label{macroscopic.FT.Btwo}
\tmq\frac{i}{6\sqrt{5}}\left( k\cdot \widehat{\mathcal{B}}-3k_2 \widehat{\mathcal{B}}_2\right)
&=-\tmTq  \partial_t\left( \tTq \widehat{\mathfrak{m}}_{10} \right)+\widehat{\mathfrak{u}}_{10},
 \\ \label{macroscopic.FT.Bjl}
\frac{1}{\sqrt{5}}  \tmq\left(ik_j \widehat{\mathcal{B}}_l+ik_l \widehat{\mathcal{B}}_j \right)
&=-\tmTq  \partial_t\left( \tTq \widehat{\mathfrak{m}}_{l+j+9} \right)+\hat{\mathfrak{u}}_{l+j+9},
\quad (l\neq j).
\end{align}
In particular using the orthonormal basis functions in \eqref{orthogonal.basis.use} we define
\begin{align} \notag
\widehat{\mathfrak{m}}_{\ell} \eqdef &~  \irnp{  \micro \hat{f}}{\tTq \mathfrak{e}_{\ell}(\tqq \pv)},
\\ \label{rhs.def.terms.FT}
\widehat{\highG}_{\ell} \eqdef &~  \irnp{\micro \hat{f}}{\tTq \frac{\pv}{| \pv|}\mathfrak{e}_{\ell}(\tqq \pv)},
\\ \notag
\hat{\mathfrak{l}}_{\ell} \eqdef &~  -\irnp{\linL(\micro \hat{f})}{\tTq \mathfrak{e}_{\ell}(\tqq \pv)},
\\ \notag
\widehat{\mathfrak{h}}_{\ell}  \eqdef &~ 
\irnp{\widehat{\Gamma}(\hat{f},\hat{f})}{\tTq \mathfrak{e}_{\ell}(\tqq \pv)},
 \\ \notag
\hat{\mathfrak{u}}_{\ell} \eqdef &~  -\tmq ik \cdot \widehat{\highG}_{\ell}
 +\hat{\mathfrak{l}}_{\ell}+\widehat{\mathfrak{h}}_{\ell},
\end{align}
where in the second to last term, from \eqref{gamma0}, we have
\begin{align}\label{def.GaF}
\widehat{\Gamma}(\hat{f},\hat{h})(k,\pv)
= &~
t^{3\qq}\int_{\mathbb{R}^3} ~  dq 
\int_{\mathbb{S}^{2}} ~ d\omega
~ v_{\Emptyset} ~  
 \sqrt{J(\tqq\qv)}
 [\hat{f}(p^{\prime})*\hat{h}(q^{\prime})](k)
 \\ \nonumber
 &~ -
 t^{3\qq}\int_{\mathbb{R}^3} ~  dq 
\int_{\mathbb{S}^{2}} ~ d\omega
~ v_{\Emptyset} ~  
 \sqrt{J(\tqq\qv)}
 [\hat{f}(p)*\hat{h}(q)](k),
\end{align}
with
\begin{eqnarray*}
{[\hat{f}(\pv^{\prime})*\hat{h}(\qv^{\prime})]}(k) & \eqdef &\int_{\Z^3_l} \hat{f}(k-l,\pv^{\prime})\hat{h}(l,\qv^{\prime})\,d\Si(l),\\
{[\hat{f}(\pv)*\hat{h}(\qv)]}(k)  & \eqdef &\int_{\Z^3_l}  \hat{f}(k-l,\pv)\hat{h}(l,\qv)\,d\Si(l).
\end{eqnarray*}
 We will also use the notation 
\begin{align}\label{index.sum.notation}
\big| \widehat{\mathfrak{m}} \big| = \sum_{\ell=1}^{13} \big| \widehat{\mathfrak{m}}_{\ell}\big|, \quad  
    \big| \widehat{\highG}\big| = \sum_{\ell=1}^{13} \big| \widehat{\highG}_{\ell}\big|, \quad 
    \big| \hat{\mathfrak{l}} \big| = \sum_{\ell=1}^{13} \big| \hat{\mathfrak{l}}_{\ell} \big|, 
    \quad 
    \big| \widehat{\mathfrak{h}} \big| = \sum_{\ell=1}^{13} \big| \widehat{\mathfrak{h}}_{\ell} \big|.
\end{align}
In the rest of these estimates we write $\widehat{\mathfrak{m}}_{\ell}$, $\widehat{\highG}_{\ell}$, $\hat{\mathfrak{l}}_{\ell}$, $\widehat{\mathfrak{h}}_{\ell}$, or $\hat{\mathfrak{u}}_{\ell}$ to denote an arbitrary term with index $\ell$ when the value of the index is unimportant. 

We take the Fourier transform of \eqref{local.cons.alt} to obtain the local conservation laws
\begin{align} \label{local.conservation.FTa}
\partial_t\widehat{\mathcal{A}}
     +
\frac{1}{2}  \tmq ik \cdot \widehat{\mathcal{B}}
+
\frac{3\qq }{t} \widehat{\mathcal{A}}
&=-\tmq ik \cdot \widehat{\highG}_{1},
\\  \label{local.conservation.FTb}
\partial_t\widehat{\mathcal{B}}_j
 +
2\tmq ik_j \left( \widehat{\mathcal{A}}
-\frac{2}{\sqrt{3}}
\widehat{\mathcal{C}}  \right)
+
\frac{3\qq }{t}\widehat{\mathcal{B}}_j
&=-\tmq ik \cdot \widehat{\highG}_{j+1}, \quad (1 \leq j \leq 3),
\\ \label{local.conservation.FTc}
\partial_t\widehat{\mathcal{C}}
  +
\frac{1}{2\sqrt{3}}  \tmq  ik \cdot \widehat{\mathcal{B}}
+
\frac{3\qq }{t} \widehat{\mathcal{C}}
&=-\tmq ik \cdot \widehat{\highG}_{5}.
\end{align}
We will use \eqref{macroscopic.FT.B}-\eqref{macroscopic.FT.Bjl} and \eqref{local.conservation.FTa}-\eqref{local.conservation.FTc} with \eqref{rhs.def.terms.FT} to prove our desired estimates.

 

First we will estimate $\widehat{\mathcal{C}}(t,k)$. We multiply \eqref{macroscopic.FT.B} by $-t^{4\qq} ik_j  \overline{\widehat{\mathcal{C}}}(t,k)$ and sum on $j$ to obtain
\begin{align*}
 \frac{2}{\sqrt{3}}\tTq |k|^2   \left|\widehat{\mathcal{C}}(t,k)\right|^2
& =   i \tq\overline{\widehat{\mathcal{C}}}(t,k)\partial_t \left(\tTq \widehat{\mathcal{B}}\cdot k +\sum_{j=1}^3 k_j   \tTq \widehat{\mathfrak{m}}_{j+13} \right)
\\
& \quad -
i t^{4\qq}\overline{\widehat{\mathcal{C}}}(t,k)  \sum_{j=1}^3 k_j\hat{\mathfrak{u}}_{j+13} .
\end{align*}
We rearrange the time derivative as
\begin{multline*}
 \frac{2}{\sqrt{3}}\tTq |k|^2   \left|\widehat{\mathcal{C}}(t,k)\right|^2
=   i \frac{d}{dt} \left( \left(\tTq \widehat{\mathcal{B}}\cdot k +\sum_{j=1}^3 k_j   \tTq \widehat{\mathfrak{m}}_{j+13} \right)\tq\overline{\widehat{\mathcal{C}}}(t,k) \right)
\\
-
i  \left( \widehat{\mathcal{B}}\cdot k +\sum_{j=1}^3 k_j    \widehat{\mathfrak{m}}_{j+13} \right) \tTq\partial_t  \left(\tq\overline{\widehat{\mathcal{C}}}(t,k)\right)
-
i t^{4\qq}\overline{\widehat{\mathcal{C}}}(t,k)  \sum_{j=1}^3 k_j\hat{\mathfrak{u}}_{j+13} .
\end{multline*}
After dividing by $|k|^2 \ge 1$ and integrating in time we have
\begin{multline*}
 \frac{2}{\sqrt{3}}\int_1^{T} \tTq|\widehat{\mathcal{C}}(t,k)|^2  dt
=
\left. \frac{i}{|k|^2} \left( \left( \widehat{\mathcal{B}}\cdot k +\sum_{j=1}^3 k_j    \widehat{\mathfrak{m}}_{j+13} \right)t^{4\qq}\overline{\widehat{\mathcal{C}}}(t,k) \right)\right|_{t=1}^T
\\
-
\frac{i}{|k|^2} \int_1^{T} \left( \widehat{\mathcal{B}}\cdot k +\sum_{j=1}^3 k_j    \widehat{\mathfrak{m}}_{j+13} \right) \tTq\partial_t  \left(\tq\overline{\widehat{\mathcal{C}}}(t,k)\right) dt
-
\frac{i}{|k|^2} \int_1^{T}t^{4\qq}\overline{\widehat{\mathcal{C}}}(t,k)  \sum_{j=1}^3 k_j\hat{\mathfrak{u}}_{j+13} dt.
\end{multline*}
We will estimate each of the terms on the right.  



From \eqref{rhs.def.terms.FT} and \eqref{dispersion.term.basis.est} we have 
\begin{align}\label{mTH.est}
\left| \widehat{\mathfrak{m}}(t,k)\right|
+
\left| \widehat{\highG}(t,k)\right|
 \lesssim 
\norm{\micro \hat{f}(t,k)}_{L^2_{\pv}}.
\end{align}
Thus, for $|k|\ge 1$ and $T \ge 1$ we have 
\begin{multline}\label{upperBIGT.est}
    \left| \frac{1}{|k|^2} \sum_{j=1}^3 
 T^{4\qq}  k_j\widehat{\mathfrak{m}}_{j+13} (T) \overline{\widehat{\mathcal{C}}}(T)\right|
 \lesssim  
 T^{4\qq}\norm{\micro \hat{f}(T,k)}_{L^2_{\pv}} \left|\widehat{\mathcal{C}}(T,k)\right|
 \\
  \lesssim  
 T^{4\qq}\Vert \hat{f}(T,k)\Vert_{L^2_{\pv}}^2
   \lesssim 
 T^{6\qq}\Vert \hat{f}(T,k)\Vert_{L^2_{\pv}}^2,
\end{multline}
and similarly 
\begin{align}\label{upperONET.est}
    \left| \frac{1}{|k|^2} \sum_{j=1}^3 
   k_j\widehat{\mathfrak{m}}_{j+13} (1) \overline{\widehat{\mathcal{C}}}(1)\right|
 \lesssim  
\norm{ \hat{f}_1(k)}_{L^2_{\pv}}^2.
\end{align}
The term  
$\left. \frac{i}{|k|^2} \left( \left( \widehat{\mathcal{B}}\cdot k  \right)t^{4\qq}\overline{\widehat{\mathcal{C}}}(t,k) \right)\right|_{t=1}^T$ also satisfies the estimates in \eqref{upperBIGT.est} and \eqref{upperONET.est} at the respective points $t=T$ and $t=1$.  Thus 
\begin{align*}
    \left|
\left. \frac{1}{|k|^2} \left( \left( \widehat{\mathcal{B}}\cdot k +\sum_{j=1}^3 k_j   \widehat{\mathfrak{m}}_{j+13} \right)t^{4\qq}\overline{\widehat{\mathcal{C}}}(t,k) \right)\right|_{t=1}^T
\right|
\lesssim
 T^{6\qq}\Vert \hat{f}(T,k)\Vert_{L^2_{\pv}}^2
 +
 \norm{ \hat{f}_1(k)}_{L^2_{\pv}}^2.
\end{align*}
Next, from \eqref{local.conservation.FTc} we have
\begin{align*}
    \partial_t\left( \tq \widehat{\mathcal{C}} \right)
  +
\frac{1}{2\sqrt{3}}   ik \cdot \widehat{\mathcal{B}}
+
\frac{2\qq }{t} \tq\widehat{\mathcal{C}}
=- ik \cdot \widehat{\highG}_{5}.
\end{align*}
Hence
\begin{align}\label{usingqleq1}
  \left|  \tTq \partial_t\left(\tq\overline{\widehat{\mathcal{C}}}\right) \right|
 \lesssim 
\tTq |k| \left( \left|\widehat{\mathcal{B}}\right|+\left| \widehat{\highG}\right| \right)
+
2\qq  t^{4\qq -1}\left|\widehat{\mathcal{C}}\right|.
\end{align}
Thus for $\qq \leq 1$ and for any small $\eta>0$ we have
\begin{multline*}
    \left|
\frac{1}{|k|^2} \int_1^{T} \left( \widehat{\mathcal{B}}\cdot k +\sum_{j=1}^3 k_j    \widehat{\mathfrak{m}}_{j+13} \right) \tTq\partial_t  \left(\tq\overline{\widehat{\mathcal{C}}}(t,k)\right) dt
\right|
\\
\leq
C_\eta\int_1^{T} \tTq\norm{\micro \hat{f}(t,k)}_{L^2_{\pv}}^2dt
+
C_\eta \int^{T}_{1} \tTq \left|\widehat{{\mathcal{B}}}(t,k)\right|^2 dt 
+
\eta \int^{T}_{1} \tTq \left|\widehat{\mathcal{C}}(t,k)\right|^2 dt. 
\end{multline*}
Now from \eqref{lin.basis.est} we obtain
\begin{align}\label{linear.op.est.grad}
\big| \widehat{\mathfrak{l}}\big|
 \lesssim 
 \tmTq
\norm{\micro \hat{f}(t,k)}_{L^2_{\pv}}.
\end{align}
With \eqref{rhs.def.terms.FT} and \eqref{linear.op.est.grad} and the previous estimates for any small $\eta>0$ we obtain
\begin{align*}
\left|  \frac{1}{|k|^2} \int_1^{T}t^{4\qq}\overline{\widehat{\mathcal{C}}}(t,k)  \sum_{j=1}^3 k_j\hat{\mathfrak{u}}_{j+13} dt
 \right|
 \leq &~
\int_1^{T}
 t^{4\qq}\left( \tmq\left|\widehat{\highG}(t,k)\right|+\big| \widehat{\mathfrak{l}}\big|+\big| \widehat{\mathfrak{h}}\big|\right) \left|\widehat{{\mathcal{C}}}(t,k)\right| dt
 \\
  \leq &~
  \eta \int^{T}_{1} \tTq  \left|\widehat{{\mathcal{C}}}(t,k)\right|^2 dt
  +C_\eta\int_1^{T} t^{5\qq}\big| \widehat{\mathfrak{h}}\big|^2dt
  \\
&~
+C_\eta\int_1^{T} \tTq\norm{\micro \hat{f}(t,k)}_{L^2_{\pv}}^2dt.
\end{align*}
We collect these estimates for $\widehat{\mathcal{C}}(t,k)$, choosing $\eta>0$ sufficiently small, to obtain
\begin{multline}\label{FT.C.estimate}
\int_1^{T} \tTq|\widehat{\mathcal{C}}(t,k)|^2  dt
\lesssim
 T^{6\qq}\Vert \hat{f}(T,k)\Vert_{L^2_{\pv}}^2+\norm{ \hat{f}_1(k)}_{L^2_{\pv}}^2
  +
\int^{T}_{1} \tTq \left|\widehat{{\mathcal{B}}}(t,k)\right|^2 dt 
 \\
+\int_1^{T} \tTq\norm{\micro \hat{f}(t,k)}_{L^2_{\pv}}^2dt
+\int_1^{T} t^{5\qq}\big| \widehat{\mathfrak{h}}\big|^2dt.
\end{multline}
This is our primary estimate for $\widehat{\mathcal{C}}(t,k)$.



Next we estimate $\widehat{\mathcal{A}}- \sqrt{3}\widehat{\mathcal{C}}$.  Multiply \eqref{macroscopic.FT.AC} by $-t^{4\qq} ik_j  \left( \overline{\widehat{\mathcal{A}}}- \sqrt{3}\overline{\widehat{\mathcal{C}}}\right)(t,k)$ and sum on $j$ to obtain
\begin{align*}
\sqrt{\frac{7}{3}}\tTq |k|^2   \left|\widehat{\mathcal{A}}(t,k)- \sqrt{3}\widehat{\mathcal{C}}(t,k)\right|^2
= &~ 
 i \tq\left( \overline{\widehat{\mathcal{A}}}- \sqrt{3}\overline{\widehat{\mathcal{C}}}\right)(t,k)\partial_t \left(\tTq\sum_{j=1}^3 k_j    \widehat{\mathfrak{m}}_{j+5} \right)
\\
&~
-
i t^{4\qq}\left( \overline{\widehat{\mathcal{A}}}- \sqrt{3}\overline{\widehat{\mathcal{C}}}\right)(t,k)  \sum_{j=1}^3 k_j\hat{\mathfrak{u}}_{j+5}.
\end{align*}
We rearrainge the time derivative to conclude that
\begin{multline*}
\sqrt{\frac{7}{3}}\tTq |k|^2   \left|\widehat{\mathcal{A}}(t,k)- \sqrt{3}\widehat{\mathcal{C}}(t,k)\right|^2
=   i \frac{d}{dt} \left(  \left(\tTq\sum_{j=1}^3 k_j    \widehat{\mathfrak{m}}_{j+5} \right)\tq\left( \overline{\widehat{\mathcal{A}}}- \sqrt{3}\overline{\widehat{\mathcal{C}}}\right) \right)
\\
-
i  \left( \sum_{j=1}^3 k_j    \widehat{\mathfrak{m}}_{j+5} \right) \tTq\partial_t  \left(\tq\left( \overline{\widehat{\mathcal{A}}}- \sqrt{3}\overline{\widehat{\mathcal{C}}}\right)\right)
-
i t^{4\qq}\left( \overline{\widehat{\mathcal{A}}}- \sqrt{3}\overline{\widehat{\mathcal{C}}}\right)  \sum_{j=1}^3 k_j\hat{\mathfrak{u}}_{j+5}.
\end{multline*}
Next divide by $|k|^2 \ge 1$ and then integrate in time to achieve
\begin{multline*}
\sqrt{\frac{7}{3}}\int_1^{T} \tTq \left|\widehat{\mathcal{A}}(t,k)- \sqrt{3}\widehat{\mathcal{C}}(t,k)\right|^2  dt
\\
=
\left. \frac{i}{|k|^2} \left(  \left(\tTq\sum_{j=1}^3 k_j    \widehat{\mathfrak{m}}_{j+5} \right)\tq\left( \overline{\widehat{\mathcal{A}}}- \sqrt{3}\overline{\widehat{\mathcal{C}}}\right)(t,k) \right)\right|_{t=1}^T
\\
-
\frac{i}{|k|^2} \int_1^{T} \left( \sum_{j=1}^3 k_j    \widehat{\mathfrak{m}}_{j+5} \right) \tTq\partial_t  \left(\tq\left( \overline{\widehat{\mathcal{A}}}- \sqrt{3}\overline{\widehat{\mathcal{C}}}\right)(t,k)\right) dt
\\
-
\frac{i}{|k|^2} \int_1^{T} t^{4\qq}\left( \overline{\widehat{\mathcal{A}}}- \sqrt{3}\overline{\widehat{\mathcal{C}}}\right)(t,k)  \sum_{j=1}^3 k_j\hat{\mathfrak{u}}_{j+5} dt.
\end{multline*}
We will estimate each of the terms on the right.  

Then as in \eqref{upperBIGT.est} and \eqref{upperONET.est} we have 
\begin{align*}
    \left|
\left. \frac{1}{|k|^2} \left(  \left(\tTq\sum_{j=1}^3 k_j    \widehat{\mathfrak{m}}_{j+5} \right)\tq\left( \overline{\widehat{\mathcal{A}}}- \sqrt{3}\overline{\widehat{\mathcal{C}}}\right)(t,k) \right)\right|_{t=1}^T
\right|
\lesssim
 T^{6\qq}\Vert \hat{f}(T,k)\Vert_{L^2_{\pv}}^2
 +
 \norm{ \hat{f}_1(k)}_{L^2_{\pv}}^2.
\end{align*}
Next subtract $\sqrt{3}$ times \eqref{local.conservation.FTc} from  \eqref{local.conservation.FTa} to obtain 
\begin{align}\notag
    \partial_t \left( \widehat{\mathcal{A}}-\sqrt{3}\widehat{\mathcal{C}} \right)
+
\frac{3\qq }{t} \left( \widehat{\mathcal{A}}-\sqrt{3}\widehat{\mathcal{C}} \right)
&=-\tmq ik \cdot \left( \widehat{\highG}_{1}-\sqrt{3}\widehat{\highG}_{5} \right).
\end{align}
In particular 
\begin{align}\notag
    \partial_t \left(\tq\left( \widehat{\mathcal{A}}-\sqrt{3}\widehat{\mathcal{C}} \right)\right)
+
\frac{2\qq }{t} \tq \left( \widehat{\mathcal{A}}-\sqrt{3}\widehat{\mathcal{C}} \right)
&=- ik \cdot \left( \widehat{\highG}_{1}-\sqrt{3}\widehat{\highG}_{5} \right).
\end{align}
Thus we have 
\begin{align*}
  \left|  \tTq \partial_t \left(\tq\left( \overline{\widehat{\mathcal{A}}}- \sqrt{3}\overline{\widehat{\mathcal{C}}}\right)(t,k)\right) \right|
 \lesssim 
\tTq |k| \left| \widehat{\highG}\right|
+
2\qq  t^{4\qq -1}\left|{\widehat{\mathcal{A}}}- \sqrt{3}{\widehat{\mathcal{C}}}\right|.
\end{align*}
We conclude, again using $\qq \leq 1$, for any small $\eta>0$ that 
\begin{multline*}
    \left|
\frac{1}{|k|^2} \int_1^{T} \left( \sum_{j=1}^3 k_j    \widehat{\mathfrak{m}}_{j+5} \right) \tTq\partial_t  \left(\tq\left( \overline{\widehat{\mathcal{A}}}- \sqrt{3}\overline{\widehat{\mathcal{C}}}\right)(t,k)\right) dt
\right|
\\
\leq
C_\eta\int_1^{T} \tTq\norm{\micro \hat{f}(t,k)}_{L^2_{\pv}}^2dt
+
\eta \int^{T}_{1} \tTq \left|{\widehat{\mathcal{A}}}- \sqrt{3}{\widehat{\mathcal{C}}}\right|^2(t,k) dt. 
\end{multline*}
Also using \eqref{rhs.def.terms.FT} and \eqref{linear.op.est.grad} and the prior estimates  for any small $\eta>0$ we obtain 
\begin{multline*}
\left|  \frac{1}{|k|^2} 
\int_1^{T} t^{4\qq}\left( \overline{\widehat{\mathcal{A}}}- \sqrt{3}\overline{\widehat{\mathcal{C}}}\right)(t,k)  \sum_{j=1}^3 k_j\hat{\mathfrak{u}}_{j+5} dt
 \right|
\\
\leq
\int_1^{T}
 t^{4\qq}\left( \tmq\left|\widehat{\highG}(t,k)\right|+\left| \widehat{\mathfrak{l}}\right|+\left| \widehat{\mathfrak{h}}\right|\right) \left|{\widehat{\mathcal{A}}}- \sqrt{3}{\widehat{\mathcal{C}}}\right|(t,k) dt
\\
\leq
C_\eta\int_1^{T} \tTq\norm{\micro \hat{f}(t,k)}_{L^2_{\pv}}^2dt
  +
  C_\eta\int_1^{T} t^{5\qq}\left| \widehat{\mathfrak{h}}\right|^2dt
+
\eta \int^{T}_{1} \tTq \left|{\widehat{\mathcal{A}}}- \sqrt{3}{\widehat{\mathcal{C}}}\right|^2(t,k) dt. 
\end{multline*}
We collect the estimates above choosing $\eta>0$ sufficiently small to obtain
\begin{multline}\label{FT.AC.estimate}
\int_1^{T} \tTq\left|\left( {\widehat{\mathcal{A}}}- \sqrt{3}{\widehat{\mathcal{C}}}\right)(t,k)\right|^2  dt
\lesssim
 T^{6\qq}\Vert \hat{f}(T,k)\Vert_{L^2_{\pv}}^2+\norm{ \hat{f}_1(k)}_{L^2_{\pv}}^2
 \\
 +
\int_1^{T} \tTq\norm{\micro \hat{f}(t,k)}_{L^2_{\pv}}^2dt
+
\int_1^{T} t^{5\qq}\left| \widehat{\mathfrak{h}}\right|^2dt.
\end{multline}
This is our primary estimate for ${\widehat{\mathcal{A}}}- \sqrt{3}{\widehat{\mathcal{C}}}$.



Now we take a linear combination of \eqref{FT.C.estimate} and \eqref{FT.AC.estimate} to obtain 
\begin{multline}\label{FT.ACt.estimate}
    \int^{T}_{1} \tTq \left|\widehat{[\mathcal{A},\mathcal{C}]}(t,k)\right|^2 dt
    \lesssim
     T^{6\qq}\Vert \hat{f}(T,k)\Vert_{L^2_{\pv}}^2+\norm{ \hat{f}_1(k)}_{L^2_{\pv}}^2
  +
 \int^{T}_{1} \tTq \left|\widehat{{\mathcal{B}}}(t,k)\right|^2 dt 
 \\
+\int_1^{T} \tTq\norm{\micro \hat{f}(t,k)}_{L^2_{\pv}}^2dt
+\int_1^{T} t^{5\qq}\left| \widehat{\mathfrak{h}}\right|^2dt.
\end{multline}
This will be our primary estimate for $\widehat{[\mathcal{A},\mathcal{C}]}(t,k)$.






Lastly we will estimate $\widehat{\mathcal{B}}$.  To this end we multiply \eqref{macroscopic.FT.Bjl} by $-ik_j \tq \overline{\widehat{\mathcal{B}}}_l$ and sum over $l\neq j$, to obtain
\begin{multline}
        -\sum_{\substack{l=1 \\ l\neq j}}^3 \left(ik_j \widehat{\mathcal{B}}_l+ik_l \widehat{\mathcal{B}}_j \right)ik_j \overline{\widehat{\mathcal{B}}}_l
    = (k_j)^2 \left(|\widehat{\mathcal{B}}(t,k)|^2 
    -
    ( \widehat{\mathcal{B}}_j )^2\right) 
    + 
     k_j \overline{\widehat{\mathcal{B}}}_j\left( k \cdot \widehat{\mathcal{B}} -k_j \widehat{\mathcal{B}}_j\right)
    \\  \label{sum.bj.calc}
     = (k_j)^2 |\widehat{\mathcal{B}}(t,k)|^2 
    + 
     k_j \overline{\widehat{\mathcal{B}}}_j\left(\left( k \cdot \widehat{\mathcal{B}} \right)
        -
2   k_j   \widehat{\mathcal{B}}_j  \right).
\end{multline}
Next, by taking suitable linear combinations of \eqref{macroscopic.FT.Bthree} and \eqref{macroscopic.FT.Btwo} we observe that for any $1\leq j\leq 3$ we have 
\begin{align}\label{macroscopic.FT.Bjjj}
    \tmq\frac{i}{6\sqrt{5}}\left( k\cdot \widehat{\mathcal{B}}-3k_j \widehat{\mathcal{B}}_j\right)
&=-\tmTq  \partial_t\left( \tTq \widehat{\mathfrak{m}}_{*j} \right)+\widehat{\mathfrak{u}}_{*j},
\end{align}
where $\widehat{\mathfrak{m}}_{*j}$ is some linear combination of  $\widehat{\mathfrak{m}}_{9}$ and $\widehat{\mathfrak{m}}_{10}$ that depends upon $j$, and  $\widehat{\mathfrak{u}}_{*j}$ is a similar linear combination of $\widehat{\mathfrak{u}}_{9}$ and $\widehat{\mathfrak{u}}_{10}$ that depends upon $j$.  Next multiply \eqref{macroscopic.FT.Bjjj} by $-ik_j \tq \overline{\widehat{\mathcal{B}}}_j$ to obtain 
\begin{align*}
    \left( k\cdot \widehat{\mathcal{B}}-2k_j \widehat{\mathcal{B}}_j\right)
    \frac{k_j \overline{\widehat{\mathcal{B}}}_j}{6\sqrt{5}}
= \frac{(k_j)^2 |\widehat{\mathcal{B}}_j|^2}{6\sqrt{5}}
+
t^{-2\qq}  ik_j\overline{\widehat{\mathcal{B}}}_j\partial_t\left( \tTq \widehat{\mathfrak{m}}_{*j} \right)
-
ik_j \tq\overline{\widehat{\mathcal{B}}}_j\widehat{\mathfrak{u}}_{*j}.
\end{align*}
Note that $|k^T\widehat{\mathcal{B}}(t,k)|^2 = \sum_{j=1}^3(k_j)^2 |\widehat{\mathcal{B}}_j|^2$.  
We plug this into \eqref{sum.bj.calc}, use \eqref{macroscopic.FT.Bjl}, and sum on $j$ to obtain 
\begin{multline*}
    |k|^2 |\widehat{\mathcal{B}}(t,k)|^2 + |k^T\widehat{\mathcal{B}}(t,k)|^2
    =
    -\sqrt{5}\sum_{j=1}^3\sum_{\substack{l=1 \\ l\neq j}}^3 \left(ik_j \widehat{\mathcal{B}}_l+ik_l \widehat{\mathcal{B}}_j \right)ik_j \overline{\widehat{\mathcal{B}}}_l
    \\
    +
6\sqrt{5}it^{-2\qq} 
\sum_{j=1}^3 
k_j\overline{\widehat{\mathcal{B}}}_j\partial_t\left( \tTq \widehat{\mathfrak{m}}_{*j} \right)
-
6\sqrt{5}i\tq
\sum_{j=1}^3
k_j \overline{\widehat{\mathcal{B}}}_j\widehat{\mathfrak{u}}_{*j}
\\
=
- t^{-2\qq}  \sqrt{5} \sum_{j=1}^3\sum_{\substack{l=1 \\ l\neq j}}^3 ik_j \overline{\widehat{\mathcal{B}}}_l\partial_t\left( \tTq \widehat{\mathfrak{m}}_{l+j+9} \right)
+
\sqrt{5}\tq\sum_{j=1}^3\sum_{\substack{l=1 \\ l\neq j}}^3 ik_j \overline{\widehat{\mathcal{B}}}_l \hat{\mathfrak{u}}_{l+j+9}
    \\
    +
6\sqrt{5}it^{-2\qq} 
\sum_{j=1}^3 
k_j\overline{\widehat{\mathcal{B}}}_j\partial_t\left( \tTq \widehat{\mathfrak{m}}_{*j} \right)
-
6\sqrt{5}i\tq
\sum_{j=1}^3
k_j \overline{\widehat{\mathcal{B}}}_j\widehat{\mathfrak{u}}_{*j}.
\end{multline*}
We have shown that 
\begin{align*}
    |k|^2 |\widehat{\mathcal{B}}(t,k)|^2 + |k^T\widehat{\mathcal{B}}(t,k)|^2
    =
t^{-2\qq} 
\sum_{l=1}^3 
k_l\overline{\widehat{\mathcal{B}}}_l\partial_t\left( \tTq \widehat{\mathfrak{m}}_{\ddagger l} \right)
-
\tq
\sum_{l=1}^3
k_l \overline{\widehat{\mathcal{B}}}_l \widehat{\mathfrak{u}}_{\ddagger l},
\end{align*}
where $\widehat{\mathfrak{m}}_{\ddagger l}$ is a linear combination of the $\widehat{\mathfrak{m}}_{l+j+9}$ and the $\widehat{\mathfrak{m}}_{*l}$, 
and $\widehat{\mathfrak{u}}_{\ddagger l}$ is a similar linear combination of the $\widehat{\mathfrak{u}}_{l+j+9}$ and the $\widehat{\mathfrak{u}}_{*l}$. 



After multiplying by $\tTq$ and dividing by $|k|^2$ we thus have
\begin{multline*}
 \tTq |\widehat{\mathcal{B}}(t,k)|^2 + \tTq \frac{|k^T\widehat{\mathcal{B}}(t,k)|^2}{|k|^2}
=
\frac{ \tq}{|k|^2}
\sum_{l=1}^3 
k_l\overline{\widehat{\mathcal{B}}}_l\partial_t\left( \tTq \widehat{\mathfrak{m}}_{\ddagger l} \right)
-
\frac{t^{4\qq} }{|k|^2}
\sum_{l=1}^3
k_l \overline{\widehat{\mathcal{B}}}_l \widehat{\mathfrak{u}}_{\ddagger l}
\\
=
\frac{d}{dt}\left(
\frac{ t^{4\qq}}{|k|^2}
\sum_{l=1}^3 
k_l\overline{\widehat{\mathcal{B}}}_l \widehat{\mathfrak{m}}_{\ddagger l} 
 \right)
 -
 \frac{1}{|k|^2 } 
 \sum_{l=1}^3 
  \tTq k_l\widehat{\mathfrak{m}}_{\ddagger l}\partial_t\left(\tq\overline{\widehat{\mathcal{B}}}^j\right) 
-
\frac{t^{4\qq} }{|k|^2}
\sum_{l=1}^3
k_l \overline{\widehat{\mathcal{B}}}_l \widehat{\mathfrak{u}}_{\ddagger l}.
\end{multline*}
Similar to the previous cases, after integrating in time over $[1,T]$, we obtain
\begin{multline*}
\int_1^{T} \tTq \left(  |\widehat{\mathcal{B}}(t,k)|^2 +  \frac{|k^T\widehat{\mathcal{B}}(t,k)|^2}{|k|^2}\right)  dt
=
 \frac{ T^{4\qq}}{|k|^2}
\sum_{l=1}^3 
k_l\overline{\widehat{\mathcal{B}}}_l(T) \widehat{\mathfrak{m}}_{\ddagger l} (T)
 \\
- \frac{ 1}{|k|^2}
\sum_{l=1}^3 
k_l\overline{\widehat{\mathcal{B}}}_l(1) \widehat{\mathfrak{m}}_{\ddagger l} (1)
  -
 \frac{1}{|k|^2 } 
 \sum_{l=1}^3 
   k_l \int_1^{T} \tTq \widehat{\mathfrak{m}}_{\ddagger l}\partial_t\left(\tq\overline{\widehat{\mathcal{B}}}_l\right)  dt
-
\frac{1 }{|k|^2}
\sum_{l=1}^3
k_l \int_1^{T}  t^{4\qq} \overline{\widehat{\mathcal{B}}}_l \widehat{\mathfrak{u}}_{\ddagger l} dt.
\end{multline*}
We will estimate each of the terms on the right side.

As in \eqref{upperBIGT.est} and \eqref{upperONET.est}, for $|k|\geq 1$, we have 
\begin{align*}
 \frac{ T^{4\qq}}{|k|^2}
\sum_{l=1}^3 
    \left| k_l\overline{\widehat{\mathcal{B}}}_l(T) \widehat{\mathfrak{m}}_{\ddagger l} (T)
\right|
+ \frac{ 1}{|k|^2}
\sum_{l=1}^3 
\left| k_l\overline{\widehat{\mathcal{B}}}_l(1) \widehat{\mathfrak{m}}_{\ddagger l} (1)\right|
\lesssim
 T^{6\qq}\Vert \hat{f}(T,k)\Vert_{L^2_{\pv}}^2
 +
 \norm{ \hat{f}_1(k)}_{L^2_{\pv}}^2.
\end{align*}
Next from \eqref{local.conservation.FTb} we have 
\begin{align*}
\tTq\partial_t\left(\tq {\widehat{\mathcal{B}}}_j\right)
 +
2\tTq ik_j \left( \widehat{\mathcal{A}}
-\frac{2}{\sqrt{3}}
\widehat{\mathcal{C}}  \right)
+
\frac{2\qq }{t}t^{4\qq}\widehat{\mathcal{B}}_j
&=- \tTq ik \cdot \widehat{\highG}_{j+1}, \quad (1 \leq j \leq 3),
\end{align*}
Hence for any $1 \leq j \leq 3$ we have
\begin{align*}
  \left|  \tTq\partial_t\left(\tq {\widehat{\mathcal{B}}}_j\right) \right|
 \lesssim 
\tTq |k| \left(\left| \widehat{\highG}\right|+\left|\widehat{[\mathcal{A},\mathcal{C}]}\right|\right)
+
2\qq  t^{4\qq -1}\left|{\widehat{\mathcal{B}}}\right|.
\end{align*}
Thus again using \eqref{mTH.est} for any $\qq \leq 1$ and for any small $\eta>0$ we have
\begin{multline*}
    \left|
 \frac{1}{|k|^2 } 
 \sum_{l=1}^3 
   k_l \int_1^{T} \tTq \widehat{\mathfrak{m}}_{\ddagger l}\partial_t\left(\tq\overline{\widehat{\mathcal{B}}}_l\right)  dt
\right|
\\
\leq
C_\eta\int_1^{T} \tTq\norm{\micro \hat{f}(t,k)}_{L^2_{\pv}}^2dt
+
\eta \int^{T}_{1} \tTq \left|\widehat{{\mathcal{B}}}(t,k)\right|^2 dt 
+
\eta \int^{T}_{1} \tTq \left|\widehat{[\mathcal{A},\mathcal{C}]}(t,k)\right|^2 dt. 
\end{multline*}
Then using \eqref{rhs.def.terms.FT}, \eqref{mTH.est}, and \eqref{linear.op.est.grad} for any small $\eta>0$ we have
\begin{align*}
\left|  \frac{1 }{|k|^2}
\sum_{l=1}^3
k_l \int_1^{T}  t^{4\qq} \overline{\widehat{\mathcal{B}}}_l \widehat{\mathfrak{u}}_{\ddagger l} dt\right|
 \leq &~
\int_1^{T}
 t^{4\qq}\left( \tmq\left|\widehat{\highG}(t,k)\right|+\left| \widehat{\mathfrak{l}}\right|+\left| \widehat{\mathfrak{h}}\right|\right) \left|\widehat{{\mathcal{B}}}(t,k)\right| dt
 \\
  \leq &~
  \eta \int^{T}_{1} \tTq  \left|\widehat{{\mathcal{B}}}(t,k)\right|^2 dt
  +C_\eta\int_1^{T} t^{5\qq}\left| \widehat{\mathfrak{h}}\right|^2dt
  \\
&~
+C_\eta\int_1^{T} \tTq\norm{\micro \hat{f}(t,k)}_{L^2_{\pv}}^2dt.
\end{align*}
We collect all of these estimates for $\widehat{\mathcal{B}}(t,k)$ to conclude for any small $\eta>0$ that 
\begin{multline}\label{FT.B.estimate}
\int_1^{T} \tTq|\widehat{\mathcal{B}}(t,k)|^2  dt
\lesssim
 T^{6\qq}\Vert \hat{f}(T,k)\Vert_{L^2_{\pv}}^2+\norm{ \hat{f}_1(k)}_{L^2_{\pv}}^2
 +
\eta \int^{T}_{1} \tTq \left|\widehat{[\mathcal{A},\mathcal{C}]}(t,k)\right|^2 dt
 \\
+C_\eta\int_1^{T} \tTq\norm{\micro \hat{f}(t,k)}_{L^2_{\pv}}^2dt
+C_\eta\int_1^{T} t^{5\qq}\big|\widehat{\mathfrak{h}}\big|^2dt.
\end{multline}
This is our primary estimate for $\widehat{\mathcal{B}}(t,k)$.



To conclude the proof we multiply \eqref{FT.ACt.estimate} by a suitably small constant $\eta_1>0$ and then we add this to \eqref{FT.B.estimate} choosing $\eta>0$ in \eqref{FT.B.estimate} small enough that $\eta-\eta_1>0$.  Combining \eqref{FT.ACt.estimate} with \eqref{FT.B.estimate} we obtain
\begin{multline*}
    \int^{T}_{1} \tTq \left|\widehat{[\mathcal{A},\mathcal{B},\mathcal{C}]}(t,k)\right|^2 dt
    \lesssim
     T^{6\qq}\Vert \hat{f}(T,k)\Vert_{L^2_{\pv}}^2+\norm{ \hat{f}_1(k)}_{L^2_{\pv}}^2
 \\
+\int_1^{T} \tTq\norm{\micro \hat{f}(t,k)}_{L^2_{\pv}}^2dt
+\int_1^{T} t^{5\qq}\left| \widehat{\mathfrak{h}}\right|^2dt.
\end{multline*}
Next we take the square root in the above inequality and use 
\begin{equation}
\label{sqrt.ineq}
\frac{1}{\sqrt{2}} (A+B)\leq \sqrt{A^2+B^2}\leq A+B,
\end{equation}
further take $\sup_{0\le t\le T}$ for the term $t^{6\qq}\Vert \hat{f}( t,k)\Vert_{L^2_{\pv}}$,
and then integrate the resulting inequality with respect to $d\Sigma(k)$ over $\Z^3$ to obtain the proof of the theorem.
\end{proof}

\end{document}
